\chapter{System model and problem formulation}

\section{Scenario}
% Celda de 500 m de radio con la BS en el centro, 15 CUE con un canal de subida cada uno, K pares V2V (5 a 30), cómo se sortean las posiciones y quiénes son vecinos (300 m o menos).

\section{Channel model}
% Las ganancias de todos los enlaces (gran y pequeña escala), la distancia mínima y qué ganancias conoce cada agente.

\section{Signal quality and rates}
% SINR de cada CUE en la BS y de cada par V2V, y las tasas de Shannon que resultan.

\section{Queue and delay}
% Dinámica de la cola de cada transmisor con llegada constante de 1 Mbps, el desborde y la relación con el retardo.

\section{Reliability}
% Probabilidad de interrupción (outage) con umbral de 5 dB y p_o = 0.1, y cómo se convierte en un umbral sobre la ganancia del enlace.

\section{The optimization problem}
% Maximizar la suma de las tasas de los CUE sujeta a tasa mínima, retardo máximo y confiabilidad de los pares, como lo plantean Li et al

\section{Formulation as a multi-agent learning problem}
% Cada par V2V es un agente; juntos forman un Dec-POMDP.

\subsection{Observation}
% El vector de 46 números que observa cada agente: su enlace en cada canal, cada CUE hacia la BS, vecinos en cada canal y su cola.

\subsection{Action}
% La pareja canal y potencia: 15 canales por 4 niveles, 60 acciones.

\subsection{Shared reward}
% La recompensa compartida que calcula la BS al final de cada slot: términos de tasa de los CUE, de colas y de confiabilidad de los pares, con sus pesos lambda.

\subsection{Objective of each agent}
% Cada agente maximiza su retorno descontado; con la recompensa compartida debe descubrir qué parte del resultado depende de su propia acción.

\subsection{What happens in a slot}
% El orden de los eventos en un slot: observar, elegir, calcular SINR, tasas y colas, recibir la recompensa y aprender.
